% Generated draft: requires mathematical and source review.
\documentclass[11pt]{article}
\usepackage[margin=25mm]{geometry}
\usepackage{fontspec}
\setmainfont{Latin Modern Roman}
\usepackage{amsmath,amssymb,mathtools,unicode-math}
\setmathfont{Latin Modern Math}
\usepackage{enumitem,needspace,xurl,hyperref}
\setlength{\parindent}{0pt}
\setlength{\parskip}{5pt}
\setlength{\emergencystretch}{3em}
\newcommand{\N}{\mathbb{N}}
\newcommand{\Z}{\mathbb{Z}}
\newcommand{\Q}{\mathbb{Q}}
\newcommand{\R}{\mathbb{R}}
\newcommand{\C}{\mathbb{C}}
\newcommand{\F}{\mathbb{F}}
\newcommand{\A}{\mathbb{A}}
\newcommand{\PP}{\mathbb{P}}
\newcommand{\E}{\mathbb{E}}
\newcommand{\eps}{\varepsilon}
\newcommand{\dd}{\,\mathrm{d}}
\begin{document}
\section*{UnsolvedMath ID: 30000985}
\textbf{Source problem code:} catalogue entry
\subsection{Definitions and mathematical statement}
Let \(p\) be a prime and \(P\) a finite \(p\)-group. The conjugacy class of \(x\in P\) is \(x^P=\{g xg^{-1}:g\in P\}\), and the class number is
\[
h_P=\#\{x^P:x\in P\}.
\]
In the customary ordinary-character interpretation, \(\operatorname{Irr}(P)\) is the set of irreducible complex characters of \(P\); \(\chi(1)\) is the degree of \(\chi\). For every integer \(k\geq0\), put
\[
a_k(P)=\#\{\chi\in\operatorname{Irr}(P):\chi(1)=p^k\}.
\]
The four requested tasks are to develop general methods that parametrize the \(P\)-conjugacy classes, compute \(h_P\), describe \(\operatorname{Irr}(P)\), and compute each \(a_k(P)\).

This is a structural and computational research program, not one stated numerical formula or a claim that finite-group data cannot be obtained by finite enumeration. The source supplies no input presentation, complexity target, preferred parametrizing objects, or distinguished family of \(p\)-groups. Its literature assessment refers to a generalized Cayley approach but gives neither the construction nor the classes to which it applies. These details are needed to formulate a specific uniform method or theorem.
\subsection{Short English statement}
Develop general structural methods for conjugacy classes, their number, irreducible characters, and the number of characters of each degree \(p^k\) in finite \(p\)-groups.
\subsection{Sources}
\begin{itemize}
\item Local UnsolvedMath registry, permanent ID \(30000985\), source code catalogue entry, ``Conjugacy Classes and Characters of Finite p-Groups.'' No individual proposer is supplied.
\item \emph{Profinite and Asymptotic Group Theory} (2009), Oberwolfach Reports, DOI 10.4171/owr/2008/28, pp. 1537--1588: \url{https://doi.org/10.4171/owr/2008/28}. This is the sole supplied external reference.
\item Supplied literature triage, checked 2026-08-21: status ``open.'' It calls this a list of general character-theory questions rather than a single assertion, and reports that the classes admitting the generalized Cayley approach are not known in the report. This dated assessment derives from supplied OWR curation; time-sensitive status and cited claims require independent verification.
\end{itemize}
\end{document}
